%!TEX program = xelatex
\documentclass{article}  
\usepackage[UTF8]{ctex}

\usepackage{amsthm,amsmath,amssymb}
\usepackage{mathrsfs}  

% 这里把间距设置为0.3 cm，你们可以自己调整
%\newcommand{\setParDis}{\setlength {\parskip} {0.1cm} }

\usepackage{setspace}
%\setlength{\parskip}{1em}


\begin{document}  
\title{无人机自组织网络}
\maketitle{}

自组网具有随时随地地自动组网，可扩展性强等特性。它是一个独立地网络自治系统，可以与基站配合协同工作。该系统能快速部署到位，建立起一套完整、强大、高抗毁的网络通信系统，提供有效的数据和多媒体通信服务。当部署多无人机进行道路交通监测工作时，可以将这些多无人机按照一定算法组织起来构建起无人机自组织网络来提升其回传数据性能。因此可以将无人机自组织网络引入至城市核心区域的网络场景中，这样形成了以地面基站提供的宏蜂窝、微蜂窝和无人机自主构建的自组织网络异构而成的空地异构无线网络。

\section{相关假设}

\begin{enumerate}
\item[(1)]无人机终端上都携带了GPS装置，以确定无人机的空间位置。
\item[(2)]每架无人机上都携带了能提供通信功能的无线信号收发设备，该设备能够实现无人机与其他无人机和基站的通信服务，为无人机回传地面交通数据构建起通信链路，为无人机之后的选网算法提供通信保障。
\item[(3)]在该场景下，无人机作为搭载数据采集设备的终端，会参考道路位置部署无人机，并且无人机需要回传数据会接入蜂窝网络。所以无人机根据道路情况会出现聚集现象，同时距离越靠近的无人机接入相同微蜂窝的概率也会更大，这样为我们的组网提供依据。
\item[(4)]无人机的组网形式采用分级结构实现。考虑将多无人机集群分成若干个簇，形成由一个簇内的簇成员节点组成的一级网络和由各簇头节点组成的二级网络构建。簇头节点分别存储了该簇的基本信息表和其他簇头节点的信息表，这样簇头节点就得到了整个自组织网络的拓扑结构。
\item[(5)]无人机的工作是持续监测道路交通数据，并不需要大范围的移动，所以分簇完成过后，簇并不会重建，簇的维护仅考虑簇内新簇头的选择。
\item[(6)]无人机在分簇之前需要指定一个全局节点以收集所有无人机的信息，在分簇结束之后也随之消失，该全局节点继续正常工作。
\item[(7)]簇类节点的相互通信使用端到端的方式完成，不需要簇头节点的转发。簇内节点要实现跨簇之间的交流需要经过簇头节点的转发。为了避免相互之间的通信干扰，簇内节点的通信与簇头节点之间的通信采用不同的频率。
\end{enumerate}

\section{自组织网络的生成与维护}
因为无人机会进行道路交通监测工作，所以其部署位置应该是在道路以及路旁上方，无人机工作位置分布会受到实际道路的限制。因此在一个城市的中心区域内，不同的无人机会被分配至不同的道路上，那么部署在同一道路上的无人机的地理位置更接近，分布的相似度更高。同时，可以借助红绿灯辅助无人机定位。\

另外，无人机监测道路交通数据的目的是为了回传这些数据给后台，那么就需要借助接入蜂窝网来传输数据，无人机之间如果关联的蜂窝网列表差异小，接入的是同一个蜂窝网，并且接收到的基站信号强度差异也很小，那么说明这些无人机的位置接近程度更高，差异更小。\

综上所述，无人机分簇的思想应该是把它们自身信息差异度小的无人机聚合在一起，换句话说，无人机之间相似度越高，它们属于同一个簇的概率更大。所以为了满足这一需求，可以采用谱聚类算法来实现无人机的分簇组网。其主要思想是把无人机看作空间中的节点，这些节点用边连接起来，“距离”较远的两个点之间的边权重值较低，而“距离”较近的两个节点之间的边权重值较高。再通过对所有节点组成的图进行切割，让切图后不同的子图之间边权重和尽可能低，而子图内的边权重和尽可能的高，从而达到聚类的目的。
\subsection{初始化图模型}
假设所有无人机飞行高度为H，用集合 $\mathcal{M} = \{1,2,\cdots,M\} $ 表示。

首先按照一级组网要求，所有无人机向全局节点发送自身信息（包括自己的三维坐标、关联的基站id及接收信号强度）。全局节点接收到之后得到所有无人机节点的信息列表，然后依次遍历节点，将每个节点与其他节点按照某种对称度量算子来计算顶点之间的相似度，从而得到所有无人机节点组成的相似矩阵 $W$ ，即：
\begin{equation*}
W = 
\begin{bmatrix}
w_{11} & w_{12} & \cdots & w_{1n} \\
w_{21} & w_{22} & \cdots & w_{2n} \\
\vdots & \vdots & \ddots & \vdots \\
w_{n1} & w_{n2} & \cdots & w_{nn} 
\end{bmatrix}
\tag{1}
\end{equation*}

$w_{ij}$就是两个节点的相似度，计算公式需要自己构想，也有一些用的比较多的高斯核函数来计算。

其次还将得到图的度矩阵D，定义为：

\begin{equation*}
D = 
\begin{bmatrix}
d_{1} &  &  &  \\
 & d_{2} &  &  \\
 &  & \ddots &  \\
 &  &  & d_{n} 
\end{bmatrix}
\tag{2}
\end{equation*}


其中对于某个顶点的度定义为 $d_{i},i = 1,2,\cdots,n$，计算公式如下。
\begin{equation}
d_{i} = \sum^{n}_{j=1}w_{ij} \tag{3}
\end{equation}
可以看出某个顶点 $v_{i}$ 的度其实就是相似矩阵第 $i$ 行的和。

\subsection{图拉普拉斯矩阵}
图拉普拉斯矩阵定义为 $L$ ：
\begin{equation}
L = D - W \tag{4}
\end{equation}
其中$D$和 $W$ 都是对称矩阵，所以L也是对称矩阵。


直接对L进行特征分解，获取其前k最小的特征值对应

\subsection{切图}





\end{document}
